% !TeX program = xelatex
\documentclass[11pt]{article}

\usepackage[a4paper,margin=1in]{geometry}
\usepackage{fontspec}
\setmainfont{texgyretermes}[
  Extension=.otf,
  UprightFont=*-regular,
  BoldFont=*-bold,
  ItalicFont=*-italic,
  BoldItalicFont=*-bolditalic
]
\usepackage{amsmath}
\usepackage[draft,demo]{graphicx}
\usepackage{booktabs}
\usepackage{csquotes}
\usepackage[
  backend=biber,
  style=authoryear,
  sorting=nyt
]{biblatex}
\setlength{\emergencystretch}{2em}

\addbibresource{references.bib}

\title{OpenLeaf Demo Document — A Local LaTeX Editor}
\author{OpenLeaf Project}
\date{\today}

\begin{document}

\maketitle

\begin{abstract}
OpenLeaf is a local editing environment for writing, compiling, and reviewing
LaTeX projects without relying on a hosted service. This demonstration document
shows a compact but realistic manuscript structure: narrative sections,
mathematics, a figure placeholder, a summary table, and a bibliography. The
source is intentionally simple so that it can serve as a clean starting point
for testing editor behavior, compiler output, and project import workflows.
\end{abstract}

\section{Purpose}

Technical writing depends on a dependable feedback loop. Authors need source
files that are easy to navigate, typography that remains consistent across
machines, and compilation feedback that makes errors actionable. A local LaTeX
editor can support that loop by keeping the project files transparent while
offering conveniences such as file browsing, live logs, PDF preview, and project
templates.

This document focuses on portable conventions. It uses XeLaTeX through
\texttt{fontspec}, selects TeX Gyre Termes as a widely available serif face, and
keeps all bibliography data in a separate BibLaTeX database. These choices make
the project suitable for a public repository while still exercising the main
features expected in a document preparation workflow.

\section{A Small Model}

LaTeX is often used because it handles structured notation and cross-referenced
material cleanly. As a simple example, suppose that an author is tracking the
quality of a writing draft over several compile-and-review cycles. Let
\(q_t\) denote a normalized quality score after cycle \(t\), and let
\(\alpha\) represent the fraction of unresolved issues corrected during a
cycle. A lightweight improvement model is
\begin{equation}
  q_{t+1} = q_t + \alpha (1 - q_t),
  \qquad 0 < \alpha \leq 1.
\end{equation}
Solving the recurrence gives the closed form
\begin{equation}
  q_t = 1 - (1 - q_0)(1 - \alpha)^t,
\end{equation}
which captures the diminishing returns of repeated revisions. The equation is
not intended as a formal theory of writing; it is included to demonstrate
display mathematics, inline notation, and the way LaTeX keeps mathematical
source readable.

\section{Document Elements}

Figure placeholders are useful in early drafts because they reserve space
before final artwork is available. The \texttt{graphicx} package is loaded in
draft mode here, so the document can show a placeholder box without requiring a
real image asset.

\begin{figure}[htbp]
  \centering
  \includegraphics[width=0.78\linewidth,height=0.28\textheight]{figures/editor-preview.pdf}
  \caption{Placeholder for an editor and preview workflow diagram.}
  \label{fig:workflow-placeholder}
\end{figure}

Table~\ref{tab:features} summarizes the main demonstration features included in
this source file. The table is intentionally modest: it verifies standard
tabular rendering without depending on custom classes or external packages
beyond common LaTeX distributions.

\begin{table}[htbp]
  \centering
  \caption{Document features exercised by this demo project.}
  \label{tab:features}
  \begin{tabular}{lll}
    \toprule
    Area & Example in this file & Purpose \\
    \midrule
    Typography & TeX Gyre Termes & Portable serif text rendering \\
    Mathematics & Display equations & Compiler and preview validation \\
    Figures & Draft graphic placeholder & Layout without binary assets \\
    Bibliography & BibLaTeX database & Citation workflow coverage \\
    \bottomrule
  \end{tabular}
\end{table}

\section{Notes on Preparation}

Good document preparation combines structure and restraint. Knuth's work on
TeX established a precise foundation for digital typesetting, while Lamport's
LaTeX system made structured authoring practical for a wide technical audience
\parencite{knuth1984texbook,lamport1994latex}. Later references on typography
and the LaTeX ecosystem remain useful because they emphasize consistency,
hierarchy, and maintainable source files
\parencite{bringhurst2012elements,mittelbach2004companion}.

For a local editor, these principles translate into practical expectations:
the source tree should be easy to inspect, compile errors should point to the
right file and line, generated artifacts should stay out of version control,
and example projects should avoid personal data. A concise demo project helps
verify those expectations without overwhelming a new user.

\printbibliography

\end{document}
